\chapter{Conclusion and Outlook}
\label{ch:conclusion}

In this chapter, we present what the OVOS implementation achieved, critically assess the limitations of the current work, and outline promising directions for future research in the area of optimized virtual orbitals and their applications in quantum computing.


\section{Conclusion}
\label{sec:con_out_conclusion}

In this thesis, we have implemented the Optimized Virtual Orbital Space (OVOS) method of Adamowicz \& Bartlett \cite{Adamowicz1987} by coding it up in Python. The following conclusions can be drawn from the results presented in Chapter~\ref{ch:results}. Below we further discuss the implications of these conclusions for quantum computing applications.


\subsection{Summary of Findings}

This section summarises the main findings of this thesis in relation to the research objectives outlined in the Introduction \ref{ch:introduction}. The implementation of the OVOS method was successfully validated against both PySCF and the original results of Adamowicz \& Bartlett, demonstrating that our code correctly captures the expected energy recovery behaviour. The key quantitative finding is that with only a fraction of the virtual orbitals (e.g., $N'_{\rm virt} \approx 50-60\% \text{ of } N_{\rm virt}$), OVOS can recover approximately 90\% of the MP2 correlation energy across all tested molecules and basis sets. Li$_2$ reached 90\% recovery at only 28\% of virtuals, matching the efficiency seen in the original work. The size of basis set affecting the number of virtual orbitals available for optimization for the same molecule, showed no clear trend in relative MP2 recovery when choosing a relative higher number of virtual orbitals in OVOS compared to the observed difference in relative lower number of virtual orbitals in OVOS for the smaller basis set. This indicates that a larger basis set, with more virtual orbitals, despite offering more flexibility for orbital optimization and does not necessarily allow for a better recovery of correlation energy, especially as we approach the full virtual space. 

The initialisation strategies for the OVOS optimization was found to mainly agree and convergence was generally robust, with cases where the number of active virtual orbitals starts approaching the full virtual space, the convergence becomes more prone to getting stuck in local minima. Here it was genereally observed across molecules that the different initialisation strategies (RHF, warm-start, random sampling) could lead to different convergence behaviours, with the random sampling and RHF, occasionally finding lower energy solutions than the warm-start. This suggests that while the OVOS method is generally robust, there may be cases where exploring multiple initialisations could yield better results, especially for larger active spaces. Thus OVOS retains the essential physics while shrinking the virtual space by a factor of two.

The OVOS optimised orbitals as VQE reference state was found to consistently reach a lower energy before the energy kink appeared, only for HF was it found that the converged energy was higher than the UHF orbitals. This suggests that the OVOS-optimised orbitals generally provide a better reference state for VQE-type algorithms compared to the default UHF orbitals, which is consistent with the theoretical expectation that OVOS should yield orbitals with better overlap with the correlated ground state. However, the fact that for HF the OVOS-optimised orbitals did not outperform the UHF orbitals indicates that there may be cases where OVOS does not provide a significant advantage, and further investigation into the factors influencing this behaviour would be valuable. It was found that UMP2 natural orbitals consistently reached a lower energy than the OVOS optimised orbitals, which is expected given that UMP2 natural orbitals are known to provide a good approximation to the correlated ground state. Only for H2O did we observe at some geometries a similar performance to the UMP2 natural orbitals. Observing the different configurations of VQE, the OVOS-optimised orbitals seems to show a bigger spread in iteration count, while the UHF and UMP2 natural orbitals seems to show a more consistent iteration count across different configurations. This suggests that while OVOS-optimised orbitals can provide a better reference state for VQE in some cases, their convergence behaviour may be more variable compared to UMP2 natural orbitals, which could have implications for the reliability of VQE calculations using OVOS orbitals. It was however observed that when regarding the mean iteration count, the OVOS-optimised orbitals seems to perform similar to the UHF and UMP2 natural orbitals, with cases where it performs better and cases where it performs worse across different configurations. So despite OVOS-optimised orbitals not beating the UMP2 natural orbitals in terms of energy convergence, it seems to perform similarly in terms of iteration count, which is an interesting result that suggests that OVOS-optimised orbitals may still be a viable option for VQE reference states, especially considering the potential qubit savings they offer.

\subsection{Critical Assessment}

Several practical limitations became apparent during implementation and testing, and they directly affect how much one can trust the present results.

The OVOS method implementation was not one to one with the original implementation by Adamowicz \& Bartlett, here we took a different approach to their initial active space. More importantly, the lack of a direct comparison with their code and the differences in computational details (e.g., integral evaluation, convergence criteria, ill-conditioning) means that while the agreement in energy recovery behaviour is reassuring, it does not guarantee that the same local minima are being found. 

Despite the overall success of the OVOS method in recovering a significant fraction of the correlation energy with a reduced virtual space, we observed some limitations and caveats. The convergence behaviour was found to be oscillatory for larger active spaces, which suggests that the optimization landscape may have multiple local minima. This could potentially be mitigated by implementing more robust optimization techniques, such as trust-region methods or line-search algorithms, rather than relying solely on the Newton-Raphson approach with level-shifting. The use of a UHF reference for closed-shell molecules introduces spin contamination, which could affect the quality of the optimized orbitals. A RHF reference would be cleaner and may lead to more stable convergence for closed-shell systems. Additionally, the block-diagonal approximation of the Hessian used in the optimization may introduce errors, especially for larger active spaces where off-diagonal coupling between different virtual orbitals could be significant.

The VQE calculations in SlowQuant were performed using a one step optimization of the wave function, which may not fully capture the potential benefits of OVOS-optimized orbitals. A more thorough exploration of the VQE optimization on quantum simulators or hardware, including multiple optimization steps and different ansatz choices, would provide a more comprehensive assessment of the advantages of OVOS orbitals for quantum algorithms. Furthermore, the comparison with UMP2 natural orbitals highlights that while OVOS can provide a better reference state than UHF orbitals, it may not always outperform more sophisticated orbital choices, and understanding the conditions under which OVOS provides the most benefit would be an important area for future investigation.

The combination of orbital-space optimisation and quantum algorithms represents a promising strategy for extending the reach of near-term quantum chemistry calculations beyond what either approach alone can achieve. This thesis provides a validated foundation for further exploration of OVOS assisted quantum computation, and the insights gained here will inform future developments in this exciting area of research.


\section{Outlook}
\label{sec:con_out_outlook}

This section outlines promising directions for future research building on the results of this thesis. The OVOS method has demonstrated significant potential for reducing the virtual orbital space while retaining most of the correlation energy, but there are many avenues for improving the algorithm, extending the method to new contexts, and applying it to larger and more complex systems. Below we discuss these opportunities in detail.


\subsection{Algorithmic Improvements}

The current OVOS implementation serves as a proof-of-concept, but there are several algorithmic improvements that could enhance its numerical stability and computational efficiency. The selection of initial active orbitals is crucial for the success of the optimization, therefore different non-trivial approaches to orbital selection could be explored. Additionally, augmenting the optimization with trust-region or line-search techniques could help prevent overshooting in regions where the rotation is large, leading to more robust convergence. Exploiting parallelism and GPU acceleration could significantly reduce the computational time, especially for larger basis sets. 

\subsection{Method Extensions}

The OVOS framework developed in this thesis is formulated at the MP2 level within an unrestricted reference, but there are several promising extensions to higher levels of theory and alternative electronic structure contexts. Extending OVOS to CCSD or CCSD(T) would allow for a more accurate description of correlation effects, and could be achieved by minimising the CCSD energy over orbital rotations. Implementing a restricted (RHF) formalism for closed-shell molecules would provide a cleaner approach that avoids UHF spin contamination and reduces the number of virtual orbitals by half. Additionally, extending OVOS to excited states, such as through EOM-CCSD or TDDFT response theory, would enable the optimisation of virtual spaces for excited-state properties. For systems with near-degeneracies, a state-averaged OVOS approach could be developed to optimise orbitals across multiple states simultaneously.


\subsection{Quantum Computing Applications}

The truncation of the virtual orbital space by OVOS is directly relevant to quantum chemistry on near-term and fault-tolerant quantum devices. The reduction in the number of virtual orbitals translates into a reduction in the number of qubits required to represent the electronic structure problem, as well as a potential reduction in circuit depth for variational algorithms like VQE. Running actual VQE calculations using OVOS orbitals on quantum simulators or hardware would provide concrete evidence of the benefits of OVOS in a quantum computing context. If OVOS orbitals indeed approximate Brueckner orbitals, we would expect a significant reduction in T1 amplitudes, which would allow us to use a UCCD ansatz instead of UCCSD, further reducing circuit depth. Quantifying the T1 reduction and benchmarking the performance of VQE with OVOS orbitals against traditional HF orbitals would be an important next step in demonstrating the practical utility of OVOS for quantum chemistry on quantum devices. 

\subsection{Larger Molecular Systems}

The OVOS method has been benchmarked on small molecules with up to 4 atoms, but extending the method to larger and more complex systems would demonstrate its practical utility for real-world applications in catalysis, drug discovery, and materials science. Transition-metal complexes, which often exhibit strong correlation effects, would be a particularly interesting test case for the limits of OVOS. Additionally, applying OVOS to biological systems such as amino acids and small peptides could provide insights into its performance in more realistic chemical environments. For periodic systems, implementing OVOS in the context of solid-state MP2 calculations would be a natural extension. Finally, combining OVOS truncation with basis-set extrapolation techniques to approach the complete basis set (CBS) limit could further enhance the accuracy of correlation energy estimates while keeping computational costs manageable.